\section{Đóng góp và bố cục của đồ án}

\subsection{Đóng góp chính}
Đồ án có năm đóng góp. Đóng góp thứ nhất là kiến trúc tách quyết định sư phạm khỏi mô hình ngôn ngữ, với sáu bất biến thiết kế được thực thi bằng cấu trúc và kiểm chứng bằng thực nghiệm. Đóng góp thứ hai là phương pháp truy xuất lai có ràng buộc theo đồ thị tri thức của chương học, giữ độ chính xác ngang truy xuất phẳng và loại kiến thức của bài sau khỏi kết quả đầu. Đóng góp thứ ba là một mô hình người học tất định, có phiên bản và phát lại được từ sổ sự kiện, cùng bộ chấm theo chuẩn viết số của Việt Nam. Đóng góp thứ tư là quy trình đưa tài liệu của giáo viên thành phương án dạy: hệ thống đọc từng trang, kiểm đáp án, dựng đồ thị tri thức, soạn phần còn thiếu và chỉ dạy bản đã được giáo viên duyệt. Đóng góp thứ năm là bộ chọn nước đi bằng contextual bandit: bộ chọn dùng LinUCB trên tập nước đi do chính sách cho phép, học từ câu trả lời kế tiếp của học sinh, có trạng thái tính lại được từ sổ sự kiện và ghi xác suất của từng lượt chọn để đánh giá ngoài chính sách.

Đồ án còn để lại một hệ thống chạy được dưới tên Monkey Edu, gồm giao diện cho học sinh, phụ huynh và giáo viên, chế độ nói chuyện bằng giọng nói và gói triển khai bằng Docker. Mã nguồn, các quyết định kiến trúc và số liệu thực nghiệm đều tái lập được từ kho mã của dự án.

\subsection{Bố cục của đồ án}
Đồ án gồm năm chương. Chương~\ref{chap:intro} giới thiệu bối cảnh, vấn đề, mục tiêu, phạm vi và đóng góp. Chương~\ref{chap:theory} trình bày cơ sở lý thuyết về hệ thống gia sư thông minh, phương pháp Socratic, mô hình ngôn ngữ lớn, truy xuất tăng cường, đồ thị tri thức và mô hình người học. Chương~\ref{chap:design} phân tích yêu cầu, nêu các bất biến thiết kế và mô tả kiến trúc cùng thiết kế chi tiết từng thành phần. Chương~\ref{chap:eval} trình bày việc triển khai, phương pháp đánh giá và kết quả thực nghiệm theo bốn câu hỏi nghiên cứu. Chương~\ref{chap:conclusion} tổng kết kết quả, nêu hạn chế và hướng phát triển.
